% ============================================================
%  Chapter 19 — Appendix: Selected Proofs
% ============================================================
\newpage
\section*{Appendix: Selected Proofs}
\addcontentsline{toc}{section}{Appendix: Selected Proofs}

This appendix collects the proofs that were deferred in the main text.
Each proof is cross-referenced to the statement it establishes.

% ---- GCD Lemma ----------------------------------------------
\begin{proofbox}[Proof --- Lemma on Coprime Divisibility (Section~\ref{sec:polynomials})]
We wish to show: if $f(x)\mid g(x)h(x)$ and $\gcd(f,g)=1$, then $f\mid h$.

Since $\gcd(f,g)=1$, by the extended Euclidean algorithm there exist
$r(x),t(x)\in\F[x]$ such that
\[
    1 = r(x)f(x) + t(x)g(x).
\]
Multiply both sides by $h(x)$:
\[
    h(x) = r(x)f(x)h(x) + t(x)g(x)h(x).
\]
Now $f\mid r(x)f(x)h(x)$ trivially, and $f\mid t(x)g(x)h(x)$ by
hypothesis ($f\mid g(x)h(x)$ implies $f\mid t(x)g(x)h(x)$).
Therefore $f\mid h(x)$. \hfill$\blacksquare$
\end{proofbox}

% ---- Intersection of subspaces is a subspace ----------------
\begin{proofbox}[Proof --- Intersection of Subspaces is a Subspace (Section~\ref{sec:span})]
Let $S,T\le V$.  We verify the three subspace conditions for $S\cap T$:
\begin{enumerate}[label=\textbf{(\roman*)}]
    \item \emph{Zero vector:} $0\in S$ and $0\in T$, so $0\in S\cap T$.
    \item \emph{Closure under addition:} If $u,v\in S\cap T$, then
          $u,v\in S$ (so $u+v\in S$) and $u,v\in T$ (so $u+v\in T$),
          hence $u+v\in S\cap T$.
    \item \emph{Closure under scalar multiplication:} If $\alpha\in\F$
          and $v\in S\cap T$, then $\alpha v\in S$ and $\alpha v\in T$,
          so $\alpha v\in S\cap T$.
\end{enumerate}
Therefore $S\cap T\le V$. \hfill$\blacksquare$
\end{proofbox}

% ---- rank A = m implies generating set ----------------------
\begin{proofbox}[Proof --- Rank Criterion for a Generating Set (Section~\ref{sec:span})]
Let $R$ be the row-echelon form of $A=[v_1\mid\cdots\mid v_n]$ obtained
by Gaussian elimination.  Then $Ax=b$ has a solution for every $b\in\F^m$
if and only if $Rx=b'$ has a solution for every $b'\in\F^m$.

If $R$ has a zero row, we can choose $b'$ so that the corresponding
equation is $0=1$, giving an inconsistency.  If $R$ has no zero rows,
every system $Rx=b'$ is consistent (each row has a pivot).  Having no zero
rows is equivalent to $\operatorname{rank}(A)=m$.  In that case $m\le n$.
\hfill$\blacksquare$
\end{proofbox}

% ---- rank A = n implies l.i. --------------------------------
\begin{proofbox}[Proof --- Rank Criterion for Linear Independence (Section~\ref{sec:lindep})]
The linear combination $\alpha_1 v_1+\cdots+\alpha_n v_n = Ax$.
The set $\{v_1,\ldots,v_n\}$ is linearly independent if and only if
$Ax=0$ has only the trivial solution $x=0$.  By the theory of row
reduction, this occurs if and only if $\operatorname{rank}(A)=n$
(every column is a pivot column), which requires $n\le m$.
\hfill$\blacksquare$
\end{proofbox}

% ---- Basis for F^m ------------------------------------------
\begin{proofbox}[Proof --- Characterisation of a Basis for $\F^m$ (Section~\ref{sec:basis})]
By the two preceding proofs:
\begin{itemize}
    \item $\{v_1,\ldots,v_n\}$ is linearly independent $\Leftrightarrow$
          $\operatorname{rank}(A)=n\le m$.
    \item $\operatorname{span}\{v_1,\ldots,v_n\}=\F^m$ $\Leftrightarrow$
          $\operatorname{rank}(A)=m\le n$.
\end{itemize}
Both conditions hold simultaneously if and only if $n=m$ and
$\operatorname{rank}(A)=n$. \hfill$\blacksquare$
\end{proofbox}

% ---- Characterisation of basis in a general f.d.v.s. --------
\begin{proofbox}[Proof --- Characterisation of Basis: $(1)\!\Leftrightarrow\!(2)\!\Leftrightarrow\!(3)$ (Section~\ref{sec:basis})]
Let $\{v_1,\ldots,v_n\}\subseteq V$ with $\dime V=n$.

\medskip\noindent
\textbf{(1)$\Rightarrow$(2):} By definition of a basis.

\medskip\noindent
\textbf{(2)$\Rightarrow$(3):} Suppose the vectors are linearly dependent.
Then there exists $j$ such that
$v_j\in\operatorname{span}\{v_i\mid i\neq j\}$, so
$\operatorname{span}\{v_1,\ldots,v_n\}\subseteq
 \operatorname{span}\{v_i\mid i\neq j\}$,
which has dimension at most $n-1<n=\dime V$, contradicting (2).

\medskip\noindent
\textbf{(3)$\Rightarrow$(1):} A linearly independent set of exactly
$n=\dime V$ vectors can be extended to a basis; but it already has the
correct size, so it is itself a basis. \hfill$\blacksquare$
\end{proofbox}

% ---- Basis subset -------------------------------------------
\begin{proofbox}[Proof --- Basis Contained in Any Spanning Set (Section~\ref{sec:basis})]
Greedily choose a maximal subset $\{w_1,\ldots,w_r\}$ of
$\{v_1,\ldots,v_n\}$ such that $w_i\notin\operatorname{span}\{w_1,\ldots,w_{i-1}\}$.
This set is linearly independent by construction.

If some $v_l\notin\operatorname{span}\{w_1,\ldots,w_r\}$, we could add $v_l$
to the set, contradicting maximality.  Hence
$\operatorname{span}\{v_1,\ldots,v_n\}\subseteq\operatorname{span}\{w_1,\ldots,w_r\}$,
so $\{w_1,\ldots,w_r\}$ is a basis for
$\operatorname{span}\{v_1,\ldots,v_n\}$. \hfill$\blacksquare$
\end{proofbox}

% ---- Existence of basis from pivot columns ------------------
\begin{proofbox}[Proof --- Pivot Columns Form a Basis (Section~\ref{sec:basis})]
Let $A=[v_1\mid\cdots\mid v_n]$ with RREF $R=[w_1\mid\cdots\mid w_n]$ and
pivot column positions $i_1<\cdots<i_r$.

\textbf{Linear independence of $\{v_{i_1},\ldots,v_{i_r}\}$:}
Suppose $x_{i_1}v_{i_1}+\cdots+x_{i_r}v_{i_r}=0$.  Define
$x\in\F^n$ with $x_{i_j}$ in position $i_j$ and $0$ elsewhere.
Then $Ax=0$, so $Rx=0$.  But $Rx=x_{i_1}e_1+\cdots+x_{i_r}e_r=0$
(since $w_{i_j}=e_j$ in RREF), giving $x_{i_j}=0$ for all $j$.

\textbf{Spanning:} For any non-pivot column $w_l$, the RREF relation
$w_l=\sum_{k=1}^r R_{kl}\,e_k$ gives, by pre-multiplying by $A$,
$v_l=\sum_{k=1}^r R_{kl}\,v_{i_k}$.
Hence $\{v_{i_1},\ldots,v_{i_r}\}$ spans $\operatorname{span}\{v_1,\ldots,v_n\}$.
\hfill$\blacksquare$
\end{proofbox}

% ---- Extension of l.i. set ----------------------------------
\begin{proofbox}[Proof --- Extending a Linearly Independent Set to a Basis (Section~\ref{sec:dimension})]
Let $\{v_1,\ldots,v_k\}$ be a linearly independent set in $V$ with
$k<n=\dime V$.  Since $V\neq\operatorname{span}\{v_1,\ldots,v_k\}$, there
exists $v_{k+1}\in V\setminus\operatorname{span}\{v_1,\ldots,v_k\}$, and
$\{v_1,\ldots,v_k,v_{k+1}\}$ is still linearly independent.  Repeating
(at most $n-k$ times) yields a linearly independent set of size $n$, which
is a basis for $V$. \hfill$\blacksquare$
\end{proofbox}

% ---- Dimension is well-defined ------------------------------
\begin{proofbox}[Proof --- Dimension is Well-Defined (Section~\ref{sec:dimension})]
Let $\{v_1,\ldots,v_m\}$ and $\{w_1,\ldots,w_n\}$ be two bases of $V$.
Every $w_j$ lies in $\operatorname{span}\{v_1,\ldots,v_m\}$.  By the
Fundamental Lemma of Linear Independence, if $n>m$ then
$\{w_1,\ldots,w_n\}$ would be linearly dependent, contradicting it being a
basis.  Hence $n\le m$.  By symmetry $m\le n$, so $m=n$. \hfill$\blacksquare$
\end{proofbox}

% ---- Dimension of a subspace --------------------------------
\begin{proofbox}[Proof --- Dimension of a Subspace is at Most $\dim V$ (Section~\ref{sec:dimension})]
Let $W\le V$ with $n=\dime V$ and $W\neq\{0\}$.  By the Fundamental
Lemma, any linearly independent subset of $W$ has at most $n$ vectors.
Build a maximal linearly independent subset $\{w_1,\ldots,w_d\}$ of $W$
by adding $w_{i+1}\notin\operatorname{span}\{w_1,\ldots,w_i\}$ at each step;
this terminates with $d\le n$.  If some $w\in W$ were outside
$\operatorname{span}\{w_1,\ldots,w_d\}$, extending would contradict
maximality; so $\{w_1,\ldots,w_d\}$ spans $W$, and $\dime W=d\le n$.
\hfill$\blacksquare$
\end{proofbox}

% ---- Sum of subspaces is a subspace -------------------------
\begin{proofbox}[Proof --- Sum of Subspaces is a Subspace (Section~\ref{sec:directsum})]
Let $W=W_1+\cdots+W_k$ with each $W_i\le V$.
\begin{enumerate}[label=\textbf{(\roman*)}]
    \item $0=0+\cdots+0\in W$ since $0\in W_i$ for all $i$.
    \item For $w=\sum w_i$ and $w'=\sum w'_i$:
          $w+w'=\sum(w_i+w'_i)\in W$ since $w_i+w'_i\in W_i$.
    \item For $\alpha\in\F$ and $w=\sum w_i$:
          $\alpha w=\sum(\alpha w_i)\in W$ since $\alpha w_i\in W_i$.
\end{enumerate}
\hfill$\blacksquare$
\end{proofbox}

% ---- Rank-Nullity --------------------------------------------
\begin{proofbox}[Proof --- Rank-Nullity Theorem (Section~\ref{sec:kernel})]
Let $A\in\F^{m\times n}$ with RREF $R$ having $r$ non-zero rows
(pivot rows).
\begin{itemize}
    \item $\operatorname{rank}(A)=r=\#\text{pivot columns}$.
    \item $\dime\Ker(A)=n-r=\#\text{free variables (non-pivot columns)}$.
\end{itemize}
Hence $\operatorname{rank}(A)+\dime\Ker(A)=r+(n-r)=n$. \hfill$\blacksquare$
\end{proofbox}

% ---- Unique linear transformation ---------------------------
\begin{proofbox}[Proof --- Existence and Uniqueness of a Linear Map Defined on a Basis (Section~\ref{sec:lt})]
\textbf{Existence.} Given a basis $\{v_1,\ldots,v_n\}$ of $V$ and
targets $w_1,\ldots,w_n\in W$, define $T:V\to W$ by
$T(\sum_{i=1}^n\alpha_i v_i)=\sum_{i=1}^n\alpha_i w_i$.  This is well-defined
since representations in a basis are unique.

\emph{Linearity:}
$T(u+v)=T(\sum(\alpha_i+\beta_i)v_i)=\sum(\alpha_i+\beta_i)w_i=T(u)+T(v)$, and
$T(\alpha v)=\alpha T(v)$.

\textbf{Uniqueness.} If $T':V\to W$ also satisfies $T'(v_i)=w_i$ for
all $i$, then for any $v=\sum\alpha_i v_i$:
$T'(v)=\sum\alpha_i T'(v_i)=\sum\alpha_i w_i=T(v)$.
Hence $T'=T$. \hfill$\blacksquare$
\end{proofbox}

% ---- Algebra of linear transformations ----------------------
\begin{proofbox}[Proof --- Algebra of Linear Transformations (Section~\ref{sec:lt})]
\textbf{Addition.} For $u,v\in V$ and $\alpha\in\F$:
\begin{align*}
    (T_1+T_2)(\alpha u+v)
    &= T_1(\alpha u+v)+T_2(\alpha u+v)\\
    &= \alpha T_1(u)+T_1(v)+\alpha T_2(u)+T_2(v)\\
    &= \alpha(T_1+T_2)(u)+(T_1+T_2)(v). \checkmark
\end{align*}

\textbf{Scalar multiplication.} For $\beta,\alpha\in\F$ and $v,v'\in V$:
\begin{align*}
    (\alpha T_1)(\beta v+v')
    &= \alpha T_1(\beta v+v')
     = \alpha\beta T_1(v)+\alpha T_1(v')
     = \beta(\alpha T_1)(v)+(\alpha T_1)(v'). \checkmark
\end{align*}
\hfill$\blacksquare$
\end{proofbox}

% ---- Composition of l.t. ------------------------------------
\begin{proofbox}[Proof --- Composition of Linear Transformations is Linear (Section~\ref{sec:lt})]
For $\alpha\in\F$ and $u,v\in V$:
\begin{align*}
    (T_2\circ T_1)(\alpha u+v)
    &= T_2(T_1(\alpha u+v))
     = T_2(\alpha T_1(u)+T_1(v))\\
    &= \alpha T_2(T_1(u))+T_2(T_1(v))
     = \alpha(T_2\circ T_1)(u)+(T_2\circ T_1)(v).
\end{align*}
\hfill$\blacksquare$
\end{proofbox}

% ---- Isomorphism preserves l.i. and basis -------------------
\begin{proofbox}[Proof --- Isomorphism Preserves Linear Independence and Bases (Section~\ref{sec:lt})]
Let $T:V\to W$ be an isomorphism and $\{v_1,\ldots,v_n\}$ a subset of
$V$.

\textbf{(i) L.i.\ is preserved:}
$\sum\alpha_i v_i=0$ iff $T(\sum\alpha_i v_i)=0$ iff
$\sum\alpha_i T(v_i)=0$ (using $T$ is injective and $T(0)=0$).

\textbf{(ii) Generating set:}
$(\Rightarrow)$ If $V=\operatorname{span}\{v_i\}$, then every $w\in W$
equals $T(v)=T(\sum\alpha_i v_i)=\sum\alpha_i T(v_i)$.
$(\Leftarrow)$ If $W=\operatorname{span}\{T(v_i)\}$, then for each $v\in V$,
$T(v)=\sum\alpha_i T(v_i)=T(\sum\alpha_i v_i)$; injectivity of $T$ gives
$v=\sum\alpha_i v_i$.

\textbf{(iii) Basis:} follows immediately from (i) and (ii). \hfill$\blacksquare$
\end{proofbox}

% ---- Isomorphic spaces have equal dimension -----------------
\begin{proofbox}[Proof --- Isomorphic Spaces Have the Same Dimension (Section~\ref{sec:lt})]
If $T:V\to W$ is an isomorphism and $\{v_1,\ldots,v_n\}$ is a basis for
$V$, then $\{T(v_1),\ldots,T(v_n)\}$ is a basis for $W$ (by the previous
proof, part (iii)).  Hence $\dime W=n=\dime V$. \hfill$\blacksquare$
\end{proofbox}

% ---- Same dimension implies isomorphic ----------------------
\begin{proofbox}[Proof --- Spaces of Equal Dimension are Isomorphic (Section~\ref{sec:lt})]
If $\dime V=\dime W=n$, then $V\cong\F^n$ (via the coordinate isomorphism)
and $\F^n\cong W$ (same reason), so $V\cong W$ by composition. \hfill$\blacksquare$
\end{proofbox}

% ---- Left-inverse equals right-inverse ----------------------
\begin{proofbox}[Proof --- Left-Inverse Equals Right-Inverse for Matrices (Section~\ref{sec:matrices})]
Suppose $BA=I$ and $AC=I$.  Then
\[
    B = BI = B(AC) = (BA)C = IC = C.
\]
\hfill$\blacksquare$
\end{proofbox}
